\section{Independent references, and the same question at two further
boundaries}
\label{sec:baselines}

\begin{figure*}[t]
\centering
\lrfigure{baselines}{\linewidth}
\caption{Per-byte cost of this data plane beside independent references. The
data plane crosses each payload byte twice, so its per-byte cost is divided by
two before being placed on the same axis.}
\label{fig:baselines}
\end{figure*}

A cost model fitted to one program's own measurements is self-consistent by
construction. Independent references are placed beside it, and
Figure~\ref{fig:baselines} shows the per-byte comparison.

\texttt{openssl speed} times the same cipher with no packet framing at all, on
the same machine, using a program this project did not write. It estimates $b$
directly and independently, and reports \refOpensslCycPerByte\ cycles per byte
against \refOursCycPerByte\ for this data plane per pass through the cipher. A
disagreement here would mean the fit was wrong; agreement is the check that it
is not.

\texttt{iperf3} over UDP measures the same interface with no cryptography and no
user-space forwarding, giving the upper bound on what any encrypted data plane
on this link could reach: \refIperfGbps\ Gb/s against \refLinerateGbps\ Gb/s for
this data plane. The third reference, a production kernel tunnel
\cite{wireguard} carrying the same traffic over the same link, is the comparison
a reader actually wants: not ``is this fast'' but ``how does its per-packet cost
compare with a deployed implementation of the same idea''. That reference did
not run on this host, so Table~\ref{tab:baselines} is absent from this build.
The two implementations would in any case differ in ways the table would have to
record, since a kernel tunnel does not cross the user-kernel boundary per packet
and this data plane decrypts and re-encrypts where a tunnel endpoint does one of
the two.

\begin{table*}[t]
\caption{Independent references. The passes column is the number of times each
payload byte crosses a cipher, which is what makes the per-byte figures
comparable.}
\label{tab:baselines}
\centering\footnotesize
\input{generated/table_baselines}
\end{table*}

\begin{figure}[t]
\centering
\lrfigure{virtualization}{\columnwidth}
\caption{Per-packet cost across isolation tiers, and along a chain of encrypting
virtual network functions.}
\label{fig:virtualization}
\end{figure}

\noindent\textbf{The isolation boundary.}
The same decomposition applies to a second boundary. A virtual network function
runs inside an isolation tier, and the expectation is that each tier adds fixed
cost per packet without adding per-byte cost, in the same shape as the system
call, one level out \cite{manousis2020}.

Four tiers are in the design: the virtual machine's own kernel with its
para-virtualised interface (\envNICDriver), a network namespace joined by a
\texttt{veth} pair, a \texttt{runc} container behind a bridge with its network
address translation, and a gVisor sandbox in which system calls are serviced by
a user-space kernel rather than by the host \cite{gvisor}. The last tier is
where $a$ should move most, since every call in the hot path then crosses into
another process.

Only \virtTiers\ tiers were measurable on this host, and the expected ordering
did not appear in them. Relative to the host tier, the namespace tier measured
\virtNetnsRelative, which is below it rather than above.
Table~\ref{tab:virtualization} gives the absolute figures. A tier that measures
cheaper than the host it sits on is not a saving from isolation. It is the
distance between two single-configuration measurements on a shared machine, and
it is the reason this report does not quote a per-tier penalty. What the two
tiers do agree on is the syscall count, which is identical across them, so the
difference is not a change in the number of boundary crossings.

\begin{table}[t]
\caption{Isolation tiers. The relative column is against the host tier, which is
the only comparison that transfers: the absolute cycle count is a property of
this machine.}
\label{tab:virtualization}
\centering\footnotesize
\input{generated/table_virtualization}
\end{table}

\noindent\textbf{This is not the cost of virtualisation.}
True bare metal is not obtainable on a public cloud, and the host tier is itself
a hypervisor guest. What is reported is the \emph{incremental} cost of each
isolation layer above the virtual machine, not the absolute cost of virtualising
a data plane. Stating that plainly matters more than the number, because a
reader who took these figures as the cost of virtualisation would understate it
by however much the hypervisor already costs, which this setup cannot see.

\noindent\textbf{Chaining.}
A service function chain composes these endpoints. Each hop terminates one
security association and originates another, which is what a chain of encrypting
functions does, and the chain is built from \texttt{veth}-joined namespaces at
lengths one through four. The slope of cost against chain length gives
\virtChainPerHop\ cycles per hop. That is the price of composition, separate
from the price of the function itself, and at this size it sits at the noise
floor of the measurement rather than above it.

\noindent\textbf{The same question at the bus.}
A GPU offload has the shape of Equation~\ref{eq:model} with different constants:
%
\begin{equation}
T(n, s) = a_{\text{gpu}} + \frac{s\,n}{B_{\text{pcie}}}
        + \frac{s\,n}{R_{\text{gpu}}},
\label{eq:gpu}
\end{equation}
%
where $a_{\text{gpu}}$ is launch and synchronisation, the middle term is the
transfer the CPU path does not pay at all, and the last is the kernel itself.
The offload wins only when the batch is large enough to amortise
$a_{\text{gpu}}$ \emph{and} the far side of the bus beats the CPU per byte. The
middle term makes the second condition much harder than it looks, because bus
bandwidth sits roughly an order of magnitude below what AES-NI sustains per
core, so for a streaming data plane the bus is usually the answer before the
kernel runs. The measurement host carried no GPU and no CUDA toolkit, so the
constants of Equation~\ref{eq:gpu} are not measured here and the dataset records
the reason.
